% ═══════════════════════════════════════════════════════════════
% LEHRERHINWEISE: Radioaktivität (Physik Klasse 10)
% Stunde 3-4 (90 Minuten)
% ═══════════════════════════════════════════════════════════════

\documentclass[11pt, a4paper]{article}

\usepackage[utf8]{inputenc}
\usepackage[T1]{fontenc}
\usepackage[ngerman]{babel}
\usepackage{helvet}
\renewcommand{\familydefault}{\sfdefault}

\usepackage[margin=2cm]{geometry}
\usepackage{enumitem}
\usepackage{tabularx}
\usepackage{booktabs}
\usepackage{array}
\usepackage{xcolor}
\usepackage{tcolorbox}
\usepackage{fancyhdr}

\setlist{nosep, leftmargin=1.5em}

% Farben
\definecolor{physik}{HTML}{2c5aa0}
\definecolor{grau}{HTML}{888888}
\definecolor{hellgrau}{HTML}{f0f0f0}
\definecolor{warnung}{HTML}{b35c00}
\definecolor{warnungbg}{HTML}{fff3e0}

% Kopfzeile
\pagestyle{fancy}
\fancyhf{}
\renewcommand{\headrulewidth}{0.4pt}
\lhead{\small\textbf{LEHRERHINWEISE} -- Physik Klasse 10}
\rhead{\small Radioaktivität (Stunde 3-4)}
\cfoot{\small\thepage}

% Boxen
\newtcolorbox{warnbox}[1][]{
    colback=warnungbg, colframe=warnung, boxrule=0.8pt, arc=0pt,
    left=6pt, right=6pt, top=6pt, bottom=6pt,
    fonttitle=\bfseries\small, title=#1
}

\newcommand{\abschnitt}[1]{%
    \vspace{10pt}\noindent\textbf{\textcolor{physik}{\large #1}}\vspace{6pt}\hrule\vspace{8pt}
}

\begin{document}

\begin{center}
{\LARGE\bfseries Lehrerhinweise: Radioaktivität}\\[6pt]
{\large Physik Klasse 10 -- Stunde 3-4 (90 Min)}\\[4pt]
{\small\textcolor{grau}{Kernphysik-Einheit, Teil 2}}
\end{center}
\vspace{8pt}\hrule\vspace{12pt}

% ═══════════════════════════════════════════════════════════════
\abschnitt{1. Stundenverlauf (90 Minuten)}

\begin{tabularx}{\textwidth}{|c|c|X|l|}
\hline
\textbf{Zeit} & \textbf{Phase} & \textbf{Inhalt} & \textbf{Material} \\
\hline
0-5 & Einstieg & Wiederholung Kernaufbau, Frage: instabile Kerne? & Tafel \\
\hline
5-20 & Input & Lernpfad: 3 Strahlungsarten (α, β, γ) & LP, Beamer \\
\hline
20-30 & Erarbeitung & AB Kasten 1: Vergleichstabelle ausfüllen & AB-02 \\
\hline
30-45 & Input & Lernpfad: Zerfallsgleichungen, Regeln & LP \\
\hline
45-55 & Erarbeitung & AB Kasten 2, Übung 2 (Gleichungen) & AB-02 \\
\hline
55-65 & Demonstration & Abschirmungsversuch mit Geiger-Zähler & Präparat, GMZ \\
\hline
65-75 & Erarbeitung & AB Seite 2: Versuchsprotokoll & AB-02 \\
\hline
75-85 & Vertiefung & Übungen 3-4 (AFB III), ionisierende Wirkung & AB-02 \\
\hline
85-90 & Sicherung & Zusammenfassung: Welche Strahlung wo gefährlich? & Tafel \\
\hline
\end{tabularx}

\vspace{12pt}

% ═══════════════════════════════════════════════════════════════
\abschnitt{2. Differenzierung}

\begin{tabularx}{\textwidth}{|c|X|X|}
\hline
\textbf{Niveau} & \textbf{Fokus} & \textbf{Hilfestellung} \\
\hline
★ (BR) & Tabelle, Eigenschaften zuordnen (Übung 1) & Farbige Merkblätter α=rot, β=blau, γ=grün \\
\hline
★★ (MR) & + Zerfallsgleichungen aufstellen & Formelkarte mit Regeln (A±?, Z±?) \\
\hline
★★★ (GY) & + AFB III: Fehler erklären, Transfer & Periodensystem zur Kontrolle \\
\hline
\end{tabularx}

\vspace{12pt}

% ═══════════════════════════════════════════════════════════════
\abschnitt{3. Häufige Fehler \& Reaktionen}

\begin{warnbox}[Typische Schülerfehler]
\begin{enumerate}
\item \textbf{„α-Strahlung ist harmlos, weil Papier sie stoppt"}\\
→ Korrektur: Bei Inkorporation (Einatmen!) ist α am gefährlichsten wegen hohem Ionisierungsvermögen.

\item \textbf{Verwechslung α-Zerfall und β-Zerfall}\\
→ Eselsbrücke: \textbf{α} = \textbf{A}bnahme um 4 (Heliumkern raus). β: Neutron → Proton.

\item \textbf{„Bei β-Zerfall werden Elektronen aus der Hülle abgestrahlt"}\\
→ Korrektur: Das Elektron entsteht IM KERN (n → p + e). Hüllenelektronen sind unbeteiligt.

\item \textbf{Massenzahl und Ordnungszahl vertauscht}\\
→ Immer: Oben = A (Masse), Unten = Z (Ladung/Protonen)
\end{enumerate}
\end{warnbox}

\vspace{12pt}

% ═══════════════════════════════════════════════════════════════
\abschnitt{4. Benötigtes Material}

\begin{itemize}
\item Lernpfad: \texttt{lernpfad-radioaktivitaet.html}
\item Arbeitsblatt: \texttt{AB-02-radioaktivitaet.pdf} (Klassensatz)
\item Musterlösung: \texttt{ML-02-radioaktivitaet.pdf}
\item \textbf{Demonstration:} Geiger-Müller-Zählrohr + Zählgerät
\item \textbf{Präparat:} Schulpräparat (z.B. Am-241 aus Rauchmelder oder Ra-226)
\item \textbf{Absorber:} Papier (1mm), Aluminium (2mm), Blei (5mm)
\item Sicherheitshinweis: Strahlenschutzverordnung beachten!
\end{itemize}

\vspace{12pt}

% ═══════════════════════════════════════════════════════════════
\abschnitt{5. Lösungen (Kurzfassung)}

\textbf{Kasten 1 (Tabelle):}
\begin{itemize}[leftmargin=2em]
\item α: Heliumkerne, +2, wenige cm Luft, Papier stoppt
\item β: Elektronen, -1, einige m Luft, Aluminium stoppt
\item γ: EM-Wellen, 0, km Reichweite, Blei stoppt (schwächt ab)
\end{itemize}

\textbf{Kasten 2:} α: A-4, Z-2 / β: A gleich, Z+1

\textbf{Übung 2:}
\begin{itemize}[leftmargin=2em]
\item a) $^{238}_{92}$U → $^{234}_{90}$Th + $^{4}_{2}$He
\item b) $^{14}_{6}$C → $^{14}_{7}$N + $^{0}_{-1}$e
\item c) $^{226}_{88}$Ra → $^{222}_{86}$Rn + $^{4}_{2}$He
\end{itemize}

\vspace{8pt}
\textit{→ Vollständige Lösungen siehe ML-02-radioaktivitaet.pdf}

\end{document}
